\documentclass[12pt,a4paper]{article}

% Drake_preamble.tex - LaTeX Preamble Template
% 作者: 謝慕揚, MD, PhD, FESC
% 
% 使用說明:
% 1. 表格請使用 longtable，不要有垂直分隔線 |
% 2. 表格內換行使用 \newline，不要使用 \\
% 3. 藥物名稱、檢驗、檢查請維持英文
% 4. Reference 格式請使用 \href 加上驗證過的 DOI

% Including essential packages for Chinese typesetting
\usepackage{xeCJK}
\usepackage{fontspec}
% Configuring Chinese font with Noto Serif CJK TC, fallback to SimSun
\IfFontExistsTF{Noto Serif CJK TC}{
  \setCJKmainfont{Noto Serif CJK TC}
  \setCJKsansfont{Noto Serif CJK TC}
  \setCJKmonofont{Noto Serif CJK TC}
}{\setCJKmainfont{SimSun}
  \setCJKsansfont{SimSun}
  \setCJKmonofont{SimSun}
}

% Including other necessary packages
\usepackage{geometry}
\usepackage{titlesec}
\usepackage{enumitem}
\usepackage{xcolor}
\usepackage{colortbl}
\usepackage{tcolorbox}
\usepackage{booktabs}
\usepackage{longtable}
\usepackage{array}
\usepackage{graphicx}
\usepackage{tikz}
\usepackage{fancyhdr}
\usepackage{multicol}
\usepackage{datetime2}
\usepackage{hyperref}
\usepackage{mdframed}
\usepackage{amsmath}
\usepackage{amssymb}
\usepackage{multirow}

\hypersetup{
    colorlinks=true,
    linkcolor=emergency_blue,
    citecolor=emergency_blue,
    urlcolor=emergency_blue,
    pdfborder={0 0 0}
}

% Defining page geometry
\geometry{
  top=2.5cm,
  bottom=2.5cm,
  left=2cm,
  right=2cm,
  headheight=25pt
}

% Adjusting topmargin to compensate for increased headheight
\addtolength{\topmargin}{-10pt}

% Defining colors
\definecolor{emergency_red}{RGB}{186,24,27}
\definecolor{emergency_blue}{RGB}{0,114,188}
\definecolor{emergency_gray}{RGB}{120,120,120}
\definecolor{highlight_yellow}{RGB}{255,250,205}

% Setting title formats
\titleformat{\section}[block]{\Large\bfseries\color{emergency_red}}{\thesection}{10pt}{}
\titleformat{\subsection}[block]{\large\bfseries\color{emergency_blue}}{\thesubsection}{8pt}{}
\titleformat{\subsubsection}[block]{\normalsize\bfseries\color{emergency_gray}}{\thesubsubsection}{6pt}{}

% Defining custom environment for important notes
\newmdenv[
    linecolor=emergency_red,
    backgroundcolor=highlight_yellow,
    linewidth=2pt,
    topline=true,
    bottomline=true,
    leftline=true,
    rightline=true,
    innertopmargin=10pt,
    innerbottommargin=10pt,
    innerleftmargin=10pt,
    innerrightmargin=10pt
]{importantbox}

% Defining note command with tcolorbox
\newcommand{\note}[1]{
    \par
    \noindent
    \begin{tcolorbox}[
        colback=emergency_blue!5!white,
        colframe=emergency_blue,
        title=\textbf{重點提醒},
        fonttitle=\bfseries,
        boxrule=0.5pt,
        arc=3pt,
        left=6pt,
        right=6pt,
        top=6pt,
        bottom=6pt
    ]
    #1
    \end{tcolorbox}
}

% Key findings box
\newcommand{\keyfinding}[1]{
    \par
    \noindent
    \begin{tcolorbox}[
        colback=yellow!10!white,
        colframe=orange!75!black,
        title=\textbf{核心發現摘要},
        fonttitle=\bfseries,
        boxrule=0.5pt,
        arc=3pt
    ]
    #1
    \end{tcolorbox}
}

% Defining custom date format for ISO 8601
\DTMnewdatestyle{iso}{
  \renewcommand{\DTMdisplaydate}[4]{\number##1-\number##2-\number##3}
  \renewcommand{\DTMDisplaydate}[4]{\number##1-\number##2-\number##3}
}
\DTMsetdatestyle{iso}
\newcommand{\mydate}{\DTMtoday}


% Custom header for this document
\pagestyle{fancy}
\fancyhf{}
\fancyhead[L]{\textcolor{emergency_red}{NEJM Clinical Decisions 教學講義}}
\fancyhead[R]{\textcolor{emergency_blue}{TAVR vs SAVR：無症狀嚴重主動脈瓣狹窄}}
\fancyfoot[C]{\textcolor{emergency_gray}{\thepage}}
\renewcommand{\headrulewidth}{1pt}
\renewcommand{\headrule}{\hbox to\headwidth{\color{emergency_red}\leaders\hrule height \headrulewidth\hfill}}

\begin{document}

%============================================================
% Title Page
%============================================================
\begin{center}
{\LARGE\bfseries\textcolor{emergency_red}{NEJM Clinical Decisions}}\\[0.5cm]
{\Large\textbf{Transcatheter or Surgical Aortic-Valve Replacement in Asymptomatic Severe Aortic Stenosis}}\\[0.3cm]
{\large 無症狀嚴重主動脈瓣狹窄：經導管或外科主動脈瓣置換術}\\[0.5cm]
{\normalsize \href{https://doi.org/10.1056/NEJMclde2416168}{N Engl J Med 2026;394:195-198. DOI: 10.1056/NEJMclde2416168}}\\[0.3cm]
{\normalsize 整理：謝慕揚 MD, PhD, FESC \quad 日期：\mydate}
\end{center}

\tableofcontents
\newpage

%============================================================
\section{前言}
%============================================================

本文為 NEJM 的「Clinical Decisions」互動式教學文章，針對一位 72 歲無症狀嚴重主動脈瓣狹窄 (aortic stenosis, AS) 合併左心室功能低下的女性病患，由兩組專家分別論述支持 TAVR 或 SAVR 的理由，供讀者參與討論並形成臨床決策。

\note{
\textbf{本文重點：}
\begin{itemize}[leftmargin=*]
    \item 無症狀嚴重 AS 但 LVEF $<$ 50\% 為瓣膜置換的 Class I 適應症
    \item 低手術風險病患中，TAVR 與 SAVR 的中期預後相似
    \item TAVR 長期（>10 年）耐久性數據仍不明
    \item 決策需由 Heart Team 評估，並與病患共同決定
\end{itemize}
}

%============================================================
\section{病例介紹 (Case Vignette)}
%============================================================

\subsection{基本資料}

\begin{longtable}{p{4cm}p{10cm}}
\toprule
\textbf{項目} & \textbf{內容} \\
\midrule
\endfirsthead
\toprule
\textbf{項目} & \textbf{內容} \\
\midrule
\endhead
\bottomrule
\endfoot
年齡/性別 & 72 歲女性 \\
\midrule
就診原因 & 例行檢查發現心雜音，轉介評估 \\
\midrule
病史 & 高血壓（使用 ramipril 2.5 mg QD）\newline 無糖尿病、無肺部疾病 \\
\midrule
生活習慣 & 不抽菸、不喝酒\newline BMI 28\newline 退休教師，每週步行 2-3 英里，三次 \\
\midrule
症狀 & \textcolor{emergency_blue}{\textbf{無症狀}}：無胸痛、無運動性呼吸困難、無疲倦、無暈厥病史 \\
\bottomrule
\end{longtable}

\subsection{身體檢查}

\begin{longtable}{p{4cm}p{10cm}}
\toprule
\textbf{生命徵象} & \textbf{數值} \\
\midrule
\endfirsthead
心率 & 75 bpm，規則 \\
\midrule
血壓 & 120/75 mmHg \\
\midrule
呼吸速率 & 16 次/分 \\
\midrule
氧飽和度 & 98\%（室內空氣） \\
\midrule
心音 & \textbf{Systolic ejection murmur}：右上胸骨旁最響亮，傳導至頸動脈 \\
\midrule
肺音 & 清晰 \\
\midrule
週邊水腫 & 無 \\
\bottomrule
\end{longtable}

\subsection{心臟超音波 (Transthoracic Echocardiography)}

\begin{longtable}{p{5cm}p{4cm}p{5cm}}
\toprule
\textbf{參數} & \textbf{數值} & \textbf{判讀} \\
\midrule
\endfirsthead
Mean pressure gradient & 42 mmHg & $\geq$ 40 mmHg = 嚴重 \\
\midrule
Aortic valve area (AVA) & 0.9 cm² & $\leq$ 1.0 cm² = 嚴重 \\
\midrule
LVEF & 44\% & \textcolor{emergency_red}{\textbf{低於正常}} \\
\bottomrule
\end{longtable}

\keyfinding{
\textbf{診斷}：\textcolor{emergency_red}{High-gradient severe aortic stenosis}
\begin{itemize}[leftmargin=*]
    \item Mean gradient $\geq$ 40 mmHg + AVA $\leq$ 1.0 cm²
    \item 無論 LVEF 或 flow 狀態如何，皆可判定為嚴重 AS
\end{itemize}
}

\subsection{進一步檢查}

\begin{itemize}[leftmargin=*]
    \item \textbf{運動負荷測試 (Stress test)}：陰性（無心肌缺血），確認無運動誘發症狀
    \item \textbf{手術風險評估}：
    \begin{itemize}
        \item STS-PROM (Society of Thoracic Surgeons Predicted Risk of Mortality)：低
        \item EuroSCORE II (European System for Cardiac Operative Risk Evaluation II)：低
    \end{itemize}
\end{itemize}

%============================================================
\section{治療決策的關鍵問題}
%============================================================

\begin{importantbox}
\textbf{核心問題}：對於這位低手術風險的無症狀嚴重 AS 病患，應建議 TAVR 還是 SAVR？
\end{importantbox}

\subsection{為何需要介入治療？}

根據 2025 ESC/EACTS Guidelines：

\begin{longtable}{p{3cm}p{8cm}p{3cm}}
\toprule
\textbf{適應症} & \textbf{說明} & \textbf{建議等級} \\
\midrule
\endfirsthead
LVEF $<$ 50\% & 排除其他原因造成的 LV dysfunction & Class I B \\
\midrule
LVEF $<$ 55\% & 可考慮早期介入 & Class IIa B \\
\bottomrule
\end{longtable}

本病患 LVEF 44\%，符合 \textbf{Class I B} 適應症，\textcolor{emergency_red}{即使無症狀也建議介入治療}。

%============================================================
\section{選項一：建議 TAVR}
%============================================================

\textit{論述者：Hélène Eltchaninoff, M.D. (法國 Rouen 大學醫院心臟科)}

\subsection{支持 TAVR 的證據}

\subsubsection{關鍵臨床試驗}

\begin{longtable}{p{3.5cm}p{5cm}p{5.5cm}}
\toprule
\textbf{試驗名稱} & \textbf{研究設計} & \textbf{主要結果} \\
\midrule
\endfirsthead
\toprule
\textbf{試驗名稱} & \textbf{研究設計} & \textbf{主要結果} \\
\midrule
\endhead
\bottomrule
\endfoot
EARLY TAVR & 多中心 RCT\newline 無症狀嚴重 AS\newline 負荷測試陰性 & TAVR vs 臨床監測\newline \textcolor{emergency_blue}{支持早期 TAVR}\newline（注意：未比較 TAVR vs SAVR） \\
\midrule
PARTNER 3 & 低風險有症狀 AS\newline 5 年追蹤 & TAVR non-inferior to SAVR \\
\midrule
Evolut Low Risk & 低風險有症狀 AS\newline 自膨式瓣膜 & TAVR non-inferior to SAVR \\
\midrule
\textcolor{emergency_red}{\textbf{RHEIA}} & \textcolor{emergency_red}{\textbf{女性專屬 RCT}}\newline 平均年齡 72 歲\newline 有症狀 AS & \textcolor{emergency_red}{\textbf{TAVR superior to SAVR}}\newline 1 年複合終點\newline（死亡、中風、再住院） \\
\bottomrule
\end{longtable}

\note{
\textbf{RHEIA 試驗的重要性}：
\begin{itemize}[leftmargin=*]
    \item 女性在 PARTNER 3 和 Evolut Low Risk 中僅佔 $<$ 40\%
    \item RHEIA 專門針對女性，顯示 balloon-expandable TAVR 優於 SAVR
    \item 差異主要來自 TAVR 組\textbf{再住院率降低}
\end{itemize}
}

\subsubsection{TAVR 的優勢}

\begin{itemize}[leftmargin=*]
    \item \textbf{微創手術}：無需開胸、無需體外循環
    \item \textbf{恢復更快}：住院時間短，可迅速恢復日常活動
    \item \textbf{適合活動型病患}：本病患每週規律運動，快速恢復對她很重要
\end{itemize}

\subsection{何時應選擇 SAVR 而非 TAVR？}

\begin{itemize}[leftmargin=*]
    \item Transfemoral TAVR 的禁忌症
    \item 解剖結構增加 TAVR 風險：
    \begin{itemize}
        \item Low coronary artery take-off（冠狀動脈開口距 annulus 過近）
        \item Shallow sinuses of Valsalva
    \end{itemize}
    \item 需要合併手術：
    \begin{itemize}
        \item 主動脈瘤
        \item 嚴重鈣化的 bicuspid valve
        \item 複雜冠狀動脈疾病
    \end{itemize}
\end{itemize}

\subsection{TAVR 論點總結}

\begin{tcolorbox}[colback=green!5!white, colframe=green!50!black, title=\textbf{支持 TAVR 的理由}]
\begin{enumerate}[leftmargin=*]
    \item 低風險病患中，中期結果 TAVR $\approx$ SAVR
    \item 女性病患有 RHEIA 試驗支持 TAVR 優於 SAVR
    \item 微創、恢復快、住院時間短
    \item 符合活動型病患的生活型態需求
    \item \textbf{重要前提}：需與病患充分溝通（shared decision making）
\end{enumerate}
\end{tcolorbox}

%============================================================
\section{選項二：建議 SAVR}
%============================================================

\textit{論述者：Elaine E. Tseng, M.D. (UCSF) 與 Vinod H. Thourani, M.D. (Piedmont Heart Institute)}

\subsection{支持 SAVR 的核心論點：長期耐久性}

\subsubsection{病患預期壽命考量}

\begin{itemize}[leftmargin=*]
    \item 美國 72 歲女性預期額外壽命：\textbf{15 年}
    \item 相同 STS-PROM 分數病患，isolated SAVR 後 8 年存活率：\textbf{85-91\%}
    \item 因此需考慮\textcolor{emergency_red}{\textbf{終身管理 (lifetime management)}}
\end{itemize}

\subsubsection{PARTNER 3 七年追蹤結果}

\begin{longtable}{p{5cm}p{4.5cm}p{4.5cm}}
\toprule
\textbf{結果指標} & \textbf{TAVR} & \textbf{SAVR} \\
\midrule
\endfirsthead
複合終點\newline（死亡、中風、再住院） & \multicolumn{2}{c}{7 年時相似（1 年時 TAVR 較優的優勢消失）} \\
\midrule
3 年後死亡率 & \textcolor{emergency_red}{持續較高} & 較低 \\
\midrule
Death or disabling stroke & HR 1.31\newline (95\% CI 0.96-1.78) & 參考組 \\
\midrule
Clinically significant valve thrombosis & \textcolor{emergency_red}{HR 5.70}\newline (95\% CI 1.29-25.25) & 參考組 \\
\bottomrule
\end{longtable}

\note{
\textbf{TAVR 長期問題的可能原因}：
\begin{itemize}[leftmargin=*]
    \item TAVR leaflet crimping（植入前瓣葉壓縮）
    \item 主動脈根部缺乏 washout jets，導致血栓形成
\end{itemize}
}

\subsubsection{瓣膜耐久性比較}

\begin{longtable}{p{4cm}p{5cm}p{5cm}}
\toprule
\textbf{項目} & \textbf{TAVR} & \textbf{SAVR} \\
\midrule
\endfirsthead
10 年以上數據 & \textcolor{emergency_red}{\textbf{目前未知}} & 已有長期數據 \\
\midrule
工程模擬預測耐久性 & \textcolor{emergency_red}{\textbf{約 8 年}} & -- \\
\midrule
$\geq$ 70 歲病患\newline 15-20 年追蹤 & -- & \textcolor{emergency_blue}{\textbf{Freedom from SVD $\geq$ 98\%}} \\
\bottomrule
\end{longtable}

\subsection{技術考量}

\begin{itemize}[leftmargin=*]
    \item \textbf{Transfemoral access 可行性}
    \item \textbf{Aortic annulus 和 aortic root 大小}
    \item \textbf{Bicuspid vs tricuspid aortic stenosis}
    \item \textbf{Coronary ostial heights}
    \item \textbf{鈣化程度}
\end{itemize}

\subsection{Patient-Prosthesis Mismatch (PPM) 的考量}

\begin{itemize}[leftmargin=*]
    \item 病患 BMI 28，需注意 PPM 風險
    \item PPM 會降低術後存活率
    \item Small annulus 現在可做 TAVR，但未來 \textbf{TAVR-in-TAVR 可能有困難}
\end{itemize}

\subsection{TAVR 失敗後的處理}

\begin{itemize}[leftmargin=*]
    \item \textbf{Surgical TAVR explantation 死亡率：12-15\%}
    \item SAVR 可保留未來 TAVR-in-SAVR 的選項
\end{itemize}

\subsection{SAVR 論點總結}

\begin{tcolorbox}[colback=blue!5!white, colframe=blue!50!black, title=\textbf{支持 SAVR 的理由}]
\begin{enumerate}[leftmargin=*]
    \item 中期結果相似，但 SAVR 長期耐久性已證實
    \item TAVR 超過 10 年的數據仍不明
    \item 工程模擬預測 TAVR 耐久性僅約 8 年
    \item SAVR 15-20 年 freedom from SVD $\geq$ 98\%
    \item 可進行 \textbf{minimally invasive SAVR + annular enlargement}
    \item 保留未來 TAVR-in-SAVR 的選項
    \item PARTNER 3 和 Evolut 試驗中 SAVR 組約 1/5 做了 concomitant procedures，可能影響短期死亡率比較
\end{enumerate}
\end{tcolorbox}

%============================================================
\section{各國指引比較}
%============================================================

\begin{longtable}{p{3cm}p{5.5cm}p{5.5cm}}
\toprule
\textbf{指引} & \textbf{建議} & \textbf{適用條件} \\
\midrule
\endfirsthead
\toprule
\textbf{指引} & \textbf{建議} & \textbf{適用條件} \\
\midrule
\endhead
\bottomrule
\endfoot
\textbf{2025 ESC/EACTS} & $\geq$ 70 歲 + tricuspid AS + 適合解剖：\textcolor{emergency_blue}{\textbf{TAVR}} & 其他 bioprosthesis 候選人：Heart Team 決定 \\
\midrule
\textbf{2020 ACC/AHA} & 65-80 歲：\newline Transfemoral TAVR 為 SAVR 的 valid alternative & 需平衡預期壽命與瓣膜耐久性\newline Shared decision making \\
\bottomrule
\end{longtable}

\note{
\textbf{指引不一致的原因}：缺乏 TAVR 在低風險、較年輕族群超過 10 年的追蹤數據。
}

%============================================================
\section{臨床決策建議}
%============================================================

\subsection{本病例的特點}

\begin{longtable}{p{4cm}p{5cm}p{5cm}}
\toprule
\textbf{特點} & \textbf{有利 TAVR} & \textbf{有利 SAVR} \\
\midrule
\endfirsthead
72 歲女性 & RHEIA 試驗支持 & 預期壽命 15 年，需考慮耐久性 \\
\midrule
活動型生活 & 快速恢復 & -- \\
\midrule
低手術風險 & 兩者皆可 & SAVR 風險可接受 \\
\midrule
LVEF 44\% & 微創減少負擔 & 需確認非缺血性原因 \\
\midrule
BMI 28 & -- & PPM 風險，SAVR 可做 annular enlargement \\
\bottomrule
\end{longtable}

\subsection{Heart Team 討論重點}

\begin{enumerate}[leftmargin=*]
    \item \textbf{解剖評估}：CT 評估 annulus、coronary heights、vascular access
    \item \textbf{Bicuspid vs Tricuspid}：文中未提及，需確認
    \item \textbf{病患意願}：充分說明兩種選項的利弊
    \item \textbf{終身管理策略}：TAVR-first vs SAVR-first approach
\end{enumerate}

%============================================================
\section{重點摘要}
%============================================================

\begin{importantbox}
\textbf{Take-home Messages:}
\begin{enumerate}[leftmargin=*]
    \item \textbf{無症狀嚴重 AS + LVEF $<$ 50\%} = Class I 瓣膜置換適應症
    \item 低風險病患中，\textbf{TAVR 與 SAVR 中期（5-7 年）預後相似}
    \item TAVR 優勢：微創、恢復快、女性有 RHEIA 支持
    \item SAVR 優勢：\textbf{長期耐久性已證實}（15-20 年）、可做 annular enlargement、保留未來選項
    \item \textbf{TAVR 超過 10 年的數據仍不明}，工程模擬預測約 8 年
    \item 決策需由 \textbf{Heart Team} 評估，並進行 \textbf{Shared Decision Making}
    \item 考量因素：年齡、預期壽命、解剖結構、病患意願、終身管理策略
\end{enumerate}
\end{importantbox}

%============================================================
\section{名詞解釋}
%============================================================

\begin{longtable}{p{4.5cm}p{9.5cm}}
\toprule
\textbf{英文術語} & \textbf{說明} \\
\midrule
\endfirsthead
\toprule
\textbf{英文術語} & \textbf{說明} \\
\midrule
\endhead
\bottomrule
\endfoot
TAVR & Transcatheter aortic-valve replacement（經導管主動脈瓣置換術） \\
\midrule
SAVR & Surgical aortic-valve replacement（外科主動脈瓣置換術） \\
\midrule
AS & Aortic stenosis（主動脈瓣狹窄） \\
\midrule
LVEF & Left ventricular ejection fraction（左心室射出分率） \\
\midrule
AVA & Aortic valve area（主動脈瓣面積） \\
\midrule
STS-PROM & Society of Thoracic Surgeons Predicted Risk of Mortality \\
\midrule
EuroSCORE II & European System for Cardiac Operative Risk Evaluation II \\
\midrule
SVD & Structural valve degeneration（瓣膜結構性退化） \\
\midrule
PPM & Patient-prosthesis mismatch（病患-瓣膜不匹配） \\
\midrule
Heart Team & 由心臟內科、心臟外科、影像科等組成的多專科團隊 \\
\midrule
TAVR-in-TAVR & 在原有 TAVR 瓣膜內再植入新的 TAVR 瓣膜 \\
\midrule
TAVR-in-SAVR & 在原有 SAVR 瓣膜內植入 TAVR 瓣膜 \\
\bottomrule
\end{longtable}

%============================================================
\section{參考文獻}
%============================================================

\begin{enumerate}
    \item Praz F, Borger MA, Lanz J, et al. 2025 ESC/EACTS guidelines for the management of valvular heart disease. \textit{Eur Heart J} 2025;46:4635-736.

    \item Généreux P, Schwartz A, Oldemeyer JB, et al. Transcatheter aortic-valve replacement for asymptomatic severe aortic stenosis. \href{https://doi.org/10.1056/NEJMoa2408527}{\textit{N Engl J Med} 2025;392:217-27.}

    \item Mack MJ, Leon MB, Thourani VH, et al. Transcatheter aortic-valve replacement in low-risk patients at five years. \href{https://doi.org/10.1056/NEJMoa2307447}{\textit{N Engl J Med} 2023;389:1949-60.}

    \item Popma JJ, Deeb GM, Yakubov SJ, et al. Transcatheter aortic-valve replacement with a self-expanding valve in low-risk patients. \href{https://doi.org/10.1056/NEJMoa1816885}{\textit{N Engl J Med} 2019;380:1706-15.}

    \item Tchetche D, Pibarot P, Bax JJ, et al. Transcatheter vs. surgical aortic valve replacement in women: the RHEIA trial. \textit{Eur Heart J} 2025;46:2079-88.

    \item Otto CM, Nishimura RA, Bonow RO, et al. 2020 ACC/AHA guideline for the management of patients with valvular heart disease. \href{https://doi.org/10.1016/j.jacc.2020.11.018}{\textit{J Am Coll Cardiol} 2021;77(4):e25-e197.}

    \item Leon MB, Mack MJ, Pibarot P, et al. Transcatheter or surgical aortic-valve replacement in low-risk patients at 7 years. \href{https://doi.org/10.1056/NEJMoa2509766}{\textit{N Engl J Med}. DOI: 10.1056/NEJMoa2509766.}

    \item Forrest JK, Yakubov SJ, Deeb GM, et al. 5-Year outcomes after transcatheter or surgical aortic valve replacement in low-risk patients with aortic stenosis. \textit{J Am Coll Cardiol} 2025;85:1523-32.

    \item Thourani VH, Habib R, Szeto WY, et al. Survival after surgical aortic valve replacement in low-risk patients: a contemporary trial benchmark. \textit{Ann Thorac Surg} 2024;117:106-12.

    \item Bourguignon T, Bouquiaux-Stablo A-L, Candolfi P, et al. Very long-term outcomes of the Carpentier-Edwards Perimount valve in aortic position. \href{https://doi.org/10.1016/j.athoracsur.2014.09.030}{\textit{Ann Thorac Surg} 2015;99:831-7.}

    \item Martin C, Sun W. Comparison of transcatheter aortic valve and surgical bioprosthetic valve durability: a fatigue simulation study. \href{https://doi.org/10.1016/j.jbiomech.2015.07.031}{\textit{J Biomech} 2015;48:3026-34.}
\end{enumerate}

\end{document}
